\section{Discussion}
\label{sec:discussion}

\subsection{Interpretation}

Our results show that LLMs create a rich landscape of unexpected partial isomorphisms between concepts from different domains.
These mappings are not random: they preferentially preserve functional roles and structural relationships, they concentrate in semantic processing layers, and they can serve as productive reasoning aids.

{\bf Why do models compress these particular concepts together?}
The most common pattern among surprising pairs involves concepts that play analogous \emph{functional roles} in their respective domains: foundation and anchor both provide stability, ingredient and enzyme both serve as transformation inputs, shield and defense both protect.
This suggests that the model's compression strategy is guided by functional abstraction: concepts that serve similar purposes get mapped to overlapping regions of representation space, regardless of whether humans group them together.
This is consistent with the Platonic Representation Hypothesis \citep{huh2024platonic}, which predicts that models converge toward representations that reflect the statistical structure of reality---and functional roles are indeed a deep structural property shared across domains.

{\bf Are these genuinely novel analogies?}
Some of our ``surprising'' pairs (\eg heart$\leftrightarrow$beat, eye$\leftrightarrow$window) are well-known human metaphors.
Their high surprise scores reflect limitations of our human similarity baseline (WordNet) rather than true novelty.
However, many top pairs (\eg ingredient$\leftrightarrow$enzyme, beam$\leftrightarrow$bearing, yield$\leftrightarrow$bearing) are not standard human metaphors.
The spontaneously generated analogies in Experiment~2C---such as \emph{chord progressions}$\leftrightarrow$\emph{market corrections} and \emph{Git}$\leftrightarrow$\emph{phylogenetic trees}---are genuinely creative and not commonly found in human discourse.

\subsection{Limitations}

\para{Human similarity baseline.}
WordNet path similarity is an imperfect proxy for human conceptual similarity.
It assigns low scores to pairs that humans might find related (\eg heart$\leftrightarrow$beat receives 0.111), inflating surprise scores for some pairs.
A more robust baseline using crowd-sourced similarity judgments or ConceptNet would strengthen the surprise metric.

\para{Sample size in Experiment 2B.}
The reasoning transfer experiment tests only 8 tasks.
While the results are consistent (unexpected hints outperform on 5 of 8 tasks), statistical power is limited.
Larger-scale evaluation with diverse reasoning benchmarks would provide stronger evidence.

\para{Self-evaluation bias.}
Using \gptfour to both generate and rate reasoning responses introduces potential bias.
The model may systematically prefer its own outputs in certain conditions.
Human evaluation or evaluation by an independent model would address this concern.

\para{Model-specific findings.}
Our results span three different models (\embmodel for embeddings, \gptfour for reasoning, \pythia for activations).
While the consistency across models strengthens the finding, each experiment is model-specific and may not generalize to all architectures.
The embedding results may not transfer to models with different training data or objectives.

\para{Causal direction.}
We demonstrate correlation between embedding surprise and activation overlap, but we do not establish that models \emph{use} shared representations for reasoning.
The shared activations could be artifacts of training data co-occurrence rather than functional compression.
Causal intervention experiments (\eg activation patching) would be needed to verify functional importance.

\subsection{Broader Implications}

\para{For interpretability.}
The partial isomorphisms we detect could serve as a lens for understanding how models organize knowledge.
The layer-specific concentration of effects (layers 5--17 in \pythia) identifies a specific locus for studying cross-domain abstraction.
Future work using sparse autoencoders \citep{bricken2023monosemanticity} could decompose polysemantic neurons to identify the specific features shared between surprising pairs.

\para{For human reasoning.}
The model-discovered analogies represent genuinely novel metaphors that could enrich human understanding.
Analogies like ingredient$\leftrightarrow$enzyme (both catalyze transformations in a recipe/reaction) or chord progressions$\leftrightarrow$market corrections (both build and release tension cyclically) could serve as pedagogical tools or creative prompts.

\para{For AI alignment.}
Understanding what concepts models ``see as similar'' reveals potential failure modes.
If a model treats shield and defense as near-synonyms, it might inappropriately transfer military reasoning to sports contexts or vice versa.
Mapping these unexpected similarities could help identify where models might make surprising generalizations.
